\section{An inverted orbit and its characters}
% BEGIN UF.Orbit
\begin{lemma}[One inverted orbit]\label{UF.Orbit}
Suppose $G=P\rtimes K$, where $P$ has odd order and
$K=\gen{a,b}$ is a Klein four group. Suppose $G=\gen{K,r}$,
where $r$ is a nonidentity involution. Write $r=ut$ with
$u\in P$, $t\in K$. Then $t\ne1$, $u^t=u^{-1}$, and
\[
 P=\gen{u,u^k}\qquad\text{for every }k\in K\setminus\gen{t}.
\]
\end{lemma}
\begin{proof}
An odd-order group has no nonidentity involution, so $t\ne1$.
The equation $(ut)^2=1$ gives $u^t=u^{-1}$. Put
$L=\gen{u,u^k}$. Since $t$ commutes with $k$, it inverts both
displayed generators of $L$. The element $k$ exchanges them. Thus
$L$ is $K$-invariant and $LK$ is a subgroup containing $K$ and $r$.
It equals $G$, and intersecting with $P$ gives $L=P$.
\end{proof}
This proof says that $t$ inverts the two indicated generators. It does
not assert that inversion is an automorphism of a nonabelian group.
% END UF.Orbit

% BEGIN UF.Characters
\begin{lemma}[Two cyclic characters]\label{UF.Characters}
In Lemma~\ref{UF.Orbit}, suppose $P$ is abelian and noncyclic.
Then it is the direct product of two nontrivial cyclic subgroups
$P_\chi$, indexed by the two characters
$\chi:K\longrightarrow\{1,-1\}$ with $\chi(t)=-1$.
The projection of $u$ generates each $P_\chi$.
\end{lemma}
\begin{proof}
Write $P$ additively. Multiplication by four is an automorphism of $P$.
The four endomorphisms
\[
 e_{\epsilon,\delta}=\tfrac14(1+\epsilon a)(1+\delta b),
 \qquad\epsilon,\delta\in\{1,-1\},
\]
are orthogonal idempotents with sum the identity: expand the products
using $a^2=b^2=1$ and $ab=ba$. Their images give the four character
components, on which $a,b$ act by $\epsilon,\delta$ respectively.
The element $t$ inverts $u$ and commutes with $K$, so it inverts the
whole $K$-orbit of $u$. That orbit generates $P$ by
Lemma~\ref{UF.Orbit}; hence $t$ acts as $-1$ on $P$.
Every component with $\chi(t)=1$ is therefore annihilated by two and
is zero. For either remaining character, the projection of each orbit
element is plus or minus $e_\chi u$. Its component is consequently
cyclic. If one component were zero then $P$ would be cyclic, so both
are nontrivial.
\end{proof}
% END UF.Characters
